\section{Multi-circuit, Reed--Muller, compound-biplane, and $A_5$ closure}
\label{sec:bt3500-3513}

The defect interval remains
\[
389\le D_{H^1}\le435,
\]
but the single-circuit method is now closed exactly. A deterministic basis exchange produces a certified 260-term circuit. Since a rank-436 fundamental circuit contains at most 436 basis coordinates while the nonzero coefficient alphabet has size $3^5-1=242$, one circuit can force at most two cancellations while adding one support. Hence no single-circuit pigeonhole proof can improve the universal upper bound below 435.

A five-channel continuous search gives a numerical candidate with ratio approximately $8.9062301931$. The exact tuple
\[
\frac1{256}(122,346,220,-240,256)
\]
has a rational-polynomial certificate proving
\[
8.905<h<9.
\]

The binary coordinate-permutation $S_3$-equivariant dimension-five code frontier was exhausted through length forty. Exact minimum points include
\[
[13,5,5],\quad[16,5,8],\quad[24,5,11],\quad[28,5,14].
\]
The sixteen columns of the $[16,5,8]$ code are the affine hyperplane
\[
\{c\in\mathbb F_2^5:\langle00111,c\rangle=1\},
\]
so the controller code is objectwise $RM(1,4)$ with weight enumerator $1+30z^8+z^{16}$.

The Clebsch biplane has thirty fourfold collision classes among its 120 two-point faults. An exhaustive search proves that three additional bits are necessary and sufficient to separate all of them. The resulting nineteen-bit readout uniquely locates all
\[
1+16+\binom{16}{2}=137
\]
fault patterns of weight at most two. The minimal companion map is cubic in the four-bit Clebsch-axis coordinate.

Finally, the Perkel and W33 rank-twenty sectors have different $A_5$ restrictions:
\[
1\oplus3\oplus3'\oplus2\cdot4\oplus5
\]
versus
\[
3\cdot1\oplus3\cdot4\oplus5.
\]
Consequently no $A_5$-equivariant rank-twenty intertwiner exists. The live chromatic boundary remains $10\le\chi(H)\le11$.
